\documentclass[11pt,openany]{book}

\usepackage[utf8]{inputenc}
\usepackage[T1]{fontenc}
\usepackage[spanish,es-tabla]{babel}
\usepackage{amsmath,amssymb,amsthm}
\usepackage{booktabs}
\usepackage{geometry}
\usepackage{hyperref}
\usepackage{listings}
\usepackage{longtable}
\usepackage{makeidx}
\usepackage{microtype}
\usepackage{url}

\geometry{letterpaper,margin=2.7cm}
\hypersetup{
  colorlinks=true,
  linkcolor=black,
  citecolor=black,
  urlcolor=blue,
  pdftitle={Manual academico de investigacion de operaciones con QSB},
  pdfauthor={Proyecto QSB}
}

\makeindex

\theoremstyle{definition}
\newtheorem{definicion}{Definicion}[chapter]
\newtheorem{ejemplo}{Ejemplo}[chapter]
\newtheorem{observacion}{Observacion}[chapter]

\newcommand{\qsb}{\textsc{QSB}}
\newcommand{\R}{\mathbb{R}}
\newcommand{\Z}{\mathbb{Z}}

\lstdefinestyle{qsbjson}{
  basicstyle=\ttfamily\small,
  breaklines=true,
  columns=fullflexible,
  frame=single,
  showstringspaces=false
}
\lstset{style=qsbjson}

\begin{document}

\frontmatter

\begin{titlepage}
  \centering
  \vspace*{2cm}
  {\Huge\bfseries Manual academico de investigacion de operaciones con \qsb\par}
  \vspace{1.5cm}
  {\Large Guia de estudio, modelacion y experimentacion computacional\par}
  \vspace{2cm}
  {\large Proyecto QSB\par}
  \vspace{0.4cm}
  {\large Editores: comunidad academica y de desarrollo de QSB\par}
  \vfill
  {\large Version viva de trabajo\par}
  {\large \today\par}
\end{titlepage}

\chapter*{Nota editorial}
\addcontentsline{toc}{chapter}{Nota editorial}
Este manual esta concebido como una publicacion academica progresiva: debe
combinar teoria establecida de investigacion de operaciones con modelos
reproducibles en \qsb. Cada capitulo debera distinguir con claridad entre
definiciones, supuestos, formulacion matematica, interpretacion gerencial,
procedimiento computacional, resultados y limitaciones.

La version inicial fija la arquitectura editorial, los estandares de citacion y
dos ejemplos basicos de programacion lineal y entera obtenidos mediante \qsb.
Las ampliaciones futuras deberian incorporar nuevos casos reales, datos de
entrada versionados, salidas verificadas y notas sobre la trazabilidad de cada
resultado.

\chapter*{Autores y contribuciones}
\addcontentsline{toc}{chapter}{Autores y contribuciones}
\begin{longtable}{p{0.28\textwidth}p{0.62\textwidth}}
\toprule
Rol & Responsabilidad editorial \\
\midrule
Proyecto QSB & Desarrollo del software, formatos normalizados y ejemplos
computacionales reproducibles. \\
Autores academicos & Curaduria teorica, demostraciones, ejercicios,
terminologia y referencias. \\
Revisores & Verificacion matematica, reproducibilidad y consistencia
pedagogica. \\
\bottomrule
\end{longtable}

\tableofcontents
\listoftables

\mainmatter

\chapter{Preliminares metodologicos}

\section{Objeto del manual}
La investigacion de operaciones estudia la construccion, analisis y solucion de
modelos cuantitativos para apoyar decisiones bajo restricciones. Este manual
adopta una perspectiva doble: presenta los fundamentos clasicos del campo y
muestra como convertirlos en modelos ejecutables con \qsb.

\section{Ciclo de modelacion}
Un estudio riguroso debe registrar, como minimo:
\begin{enumerate}
  \item el sistema real y la decision que se desea apoyar;
  \item las variables de decision, parametros y unidades;
  \item los supuestos que vuelven tratable el modelo;
  \item la funcion objetivo y las restricciones;
  \item el procedimiento de solucion;
  \item la interpretacion, validacion y sensibilidad del resultado.
\end{enumerate}

\section{Trazabilidad computacional en QSB}
Los ejemplos del manual usaran archivos de entrada en
\texttt{manual/examples/} y salidas obtenidas con comandos de \qsb. Cuando un
ejemplo provenga de datos
externos, debera documentarse su fuente, licencia, fecha de extraccion,
transformaciones y unidad de analisis.

\begin{observacion}
Un resultado computacional no sustituye la interpretacion. El solucionador
produce valores optimos bajo los supuestos codificados; el analista debe decidir
si esos supuestos son defendibles para el sistema real.
\end{observacion}

\chapter{Programacion lineal}

\section{Modelo canonico}
Un problema de programacion lineal (PL) optimiza una funcion lineal sobre una
region factible descrita por igualdades o desigualdades lineales
\cite{dantzig1963linear,bertsimas1997introduction}.

\begin{definicion}
Sea \(x \in \R^n\) el vector de variables de decision, \(c \in \R^n\), \(A \in
\R^{m \times n}\) y \(b \in \R^m\). Una forma estandar de maximizacion es
\[
  \max \; c^\top x
  \quad \text{sujeto a} \quad
  Ax \leq b,\; x \geq 0.
\]
\end{definicion}

\section{Ejemplo QSB: mezcla de produccion}
Considere dos productos \(X_1\) y \(X_2\). Cada unidad de \(X_1\) aporta 50
unidades monetarias y cada unidad de \(X_2\) aporta 60. Dos recursos limitan la
produccion:
\[
\begin{aligned}
  \max \quad & 50X_1 + 60X_2 \\
  \text{s.a.}\quad & 2X_1 + 3X_2 \leq 180,\\
  & 3X_1 + 2X_2 \leq 150,\\
  & X_1, X_2 \geq 0.
\end{aligned}
\]

El archivo normalizado usado por \qsb{} esta en
\texttt{manual/examples/lp\_sample.json}. La solucion fue obtenida con:

\begin{lstlisting}
swift run qsb solve-json manual/examples/lp_sample.json
\end{lstlisting}

\begin{table}[h]
\centering
\caption{Solucion del ejemplo de programacion lineal}
\begin{tabular}{lr}
\toprule
Cantidad & Valor \\
\midrule
\(X_1\) & 18 \\
\(X_2\) & 48 \\
Valor objetivo & 3780 \\
\bottomrule
\end{tabular}
\end{table}

\section{Interpretacion}
La solucion usa completamente ambos recursos:
\[
2(18)+3(48)=180,\qquad 3(18)+2(48)=150.
\]
Por tanto, las dos restricciones son activas. En una version ampliada del
manual, este ejemplo deberia complementarse con precios sombra, rangos de
sensibilidad y una explicacion geometrica de los vertices factibles.

\chapter{Programacion entera}

\section{Motivacion}
La programacion entera extiende la programacion lineal al exigir que algunas o
todas las variables adopten valores enteros. Esta restriccion permite modelar
decisiones indivisibles como abrir instalaciones, seleccionar proyectos,
asignar turnos o producir lotes discretos \cite{nemhauser1988integer,wolsey1998integer}.

\section{Modelo general}
Un problema mixto entero lineal puede escribirse como
\[
  \min \; c^\top x + d^\top y
  \quad \text{sujeto a} \quad
  Ax + By \geq b,\;
  x \in \R_+^n,\;
  y \in \Z_+^p.
\]
La dificultad computacional crece porque la region factible deja de ser
convexa en el sentido continuo: la relajacion lineal puede proponer soluciones
fraccionarias que no son implementables.

\section{Ejemplo QSB: requerimientos minimos discretos}
El ejemplo \texttt{manual/examples/ilp\_sample.json} minimiza el costo de dos
actividades enteras:
\[
\begin{aligned}
  \min \quad & 2.5X_1 + 2X_2 \\
  \text{s.a.}\quad & 6X_1 + 3X_2 \geq 200,\\
  & 3X_1 + 5X_2 \geq 180,\\
  & X_1, X_2 \in \Z_+.
\end{aligned}
\]

Comando reproducible:
\begin{lstlisting}
swift run qsb solve-json-ilp manual/examples/ilp_sample.json
\end{lstlisting}

\begin{table}[h]
\centering
\caption{Solucion del ejemplo de programacion entera}
\begin{tabular}{lr}
\toprule
Cantidad & Valor \\
\midrule
\(X_1\) & 22 \\
\(X_2\) & 23 \\
Valor objetivo & 101 \\
\bottomrule
\end{tabular}
\end{table}

\section{Notas de verificacion}
La solucion satisface las restricciones:
\[
6(22)+3(23)=201 \geq 200,\qquad 3(22)+5(23)=181 \geq 180.
\]
La cercania de ambas restricciones a sus requerimientos muestra el equilibrio
tipico entre costo minimo e integridad de las decisiones.

\chapter{Modelos de redes}

\section{Estructura comun}
Un modelo de redes representa entidades como nodos y relaciones como arcos. Esta
abstraccion cubre rutas mas cortas, arboles de expansion minima, flujo maximo,
transporte, asignacion y vendedor viajero \cite{ahuja1993network}.

\section{Familias contempladas en QSB}
El formato normalizado de \qsb{} usa un sobre JSON con un campo
\texttt{kind} y un objeto \texttt{model}. Las familias documentadas actualmente
incluyen:
\begin{longtable}{ll}
\toprule
Codigo & Familia \\
\midrule
SPP & Ruta mas corta \\
MST & Arbol de expansion minima \\
MFP & Flujo maximo \\
TSP & Vendedor viajero \\
AP & Asignacion \\
TP & Transporte \\
\bottomrule
\end{longtable}

\section{Criterio editorial para nuevos casos}
Cada caso de red debera incluir diagrama, matriz o lista de arcos, formulacion
matematica, archivo JSON de entrada, comando de solucion y verificacion manual
del costo, flujo o asignacion resultante.

\chapter{Modelos de inventarios}

\section{Decision fundamental}
Los modelos de inventarios equilibran costos de ordenar, mantener, agotar y
reponer existencias. El problema central es decidir cuanto pedir o producir, y
cuando hacerlo, bajo una demanda determinista o aleatoria
\cite{harris1913many,silver1998inventory}.

\section{Modelo EOQ}
En el modelo clasico de cantidad economica de pedido, con demanda anual \(D\),
costo fijo por orden \(S\) y costo anual de mantener una unidad \(H\), la
cantidad optima es
\[
  Q^* = \sqrt{\frac{2DS}{H}}.
\]

\section{Uso futuro con QSB}
El manual debe incorporar ejemplos reproducibles para EOQ, EOQ con descuentos,
modelo del vendedor de periodicos y dimensionamiento de lotes cuando los casos
esten disponibles como entradas versionadas y salidas verificadas.

\chapter{Modelos de colas}

\section{Proposito}
La teoria de colas estudia sistemas donde clientes, trabajos o unidades esperan
servicio ante capacidad limitada. Sus metricas centrales incluyen utilizacion,
longitud esperada de cola, tiempo esperado en el sistema y probabilidad de
bloqueo \cite{gross2008fundamentals}.

\section{Modelo M/M/1}
Para llegadas Poisson con tasa \(\lambda\), tiempos de servicio exponenciales
con tasa \(\mu\) y un servidor, si \(\rho=\lambda/\mu<1\), entonces
\[
  L = \frac{\rho}{1-\rho}, \qquad
  L_q = \frac{\rho^2}{1-\rho}, \qquad
  W = \frac{1}{\mu-\lambda}, \qquad
  W_q = \frac{\lambda}{\mu(\mu-\lambda)}.
\]

\section{Criterio de documentacion}
Los ejemplos de colas en \qsb{} deberan registrar la notacion Kendall, la unidad
de tiempo, las distribuciones asumidas, la tasa efectiva de llegada y las notas
del backend cuando se use una aproximacion.

\chapter{Pronosticos}

\section{Rol en investigacion de operaciones}
Los pronosticos alimentan modelos de inventarios, capacidad, programacion y
planeacion agregada. Una prediccion debe evaluarse por su utilidad operacional,
no solo por ajuste historico \cite{hyndman2021forecasting}.

\section{Metodos iniciales}
El manual cubrira, de forma progresiva, promedio movil, suavizamiento
exponencial, estacionalidad y regresion. Para cada metodo se debe reportar:
serie original, transformaciones, horizonte, parametros, errores y criterio de
seleccion.

\section{Conexion con QSB}
Los comandos de \qsb{} para series de tiempo permiten fijar ventanas,
coeficientes de suavizamiento y horizontes. Cada ejemplo publicado debera
conservar el archivo de datos y el comando exacto usado.

\chapter{Calendarizacion y secuenciacion}

\section{Definicion}
La calendarizacion asigna trabajos a recursos en el tiempo. En manufactura y
servicios, sus objetivos frecuentes son minimizar el makespan, tardanza,
tiempos ociosos o costos por incumplimiento \cite{pinedo2016scheduling}.

\section{Flow shop y job shop}
En un flow shop, todos los trabajos visitan las maquinas en el mismo orden. En
un job shop, cada trabajo puede tener una ruta propia. Esta diferencia altera
profundamente la dificultad del problema y el tipo de representacion requerida.

\section{Estandar para ejemplos}
Los ejemplos del manual deberan incluir tabla de operaciones, diagrama de Gantt,
makespan, secuencia de trabajos y salida JSON de \qsb{} con tiempos de inicio y
fin por operacion.

\chapter{Referencia de terminos}

\begin{description}
  \item[Funcion objetivo] Expresion matematica que cuantifica el criterio que
  se desea maximizar o minimizar.
  \item[Variable de decision] Magnitud controlable por el analista o por el
  decisor dentro del modelo.
  \item[Parametro] Dato exogeno que el modelo toma como conocido durante la
  solucion.
  \item[Restriccion] Condicion que limita los valores admisibles de las
  variables de decision.
  \item[Region factible] Conjunto de soluciones que satisfacen todas las
  restricciones.
  \item[Solucion optima] Solucion factible que alcanza el mejor valor de la
  funcion objetivo.
  \item[Relajacion lineal] Modelo obtenido al reemplazar restricciones de
  integralidad por variables continuas.
  \item[Utilizacion] Fraccion de la capacidad de servicio empleada en promedio.
  \item[Makespan] Tiempo de finalizacion del ultimo trabajo en un programa de
  produccion o servicio.
\end{description}


\backmatter
\bibliographystyle{apalike}
\bibliography{references}
\printindex

\end{document}
